\section{Related Work}\label{sec:related}

\subsection{Formalized Complexity Theory and Mechanized Reductions}

Machine verification of complexity-theoretic arguments remains substantially less mature than formal verification in algebra, analysis, or standard algorithmics. Forster et al.~\cite{forster2019verified} developed certified machine models and computability infrastructure in Coq, and Kunze et al.~\cite{kunze2019formal} formalized major proof-theoretic components in Coq as well. In Isabelle/HOL, much of the mature work has centered on verified algorithms and resource analysis for concrete procedures rather than on families of hardness reductions~\cite{nipkow2002isabelle,haslbeck2021verified}. Lean 4 and Mathlib provide increasingly capable foundations for mechanized finite mathematics and computation~\cite{mathlib2020,moura2021lean4}, but reusable reduction suites for coNP/$\Sigma_2^P$/PP/PSPACE-style arguments remain relatively rare.

The present artifact is intended as a problem-specific certification layer within that broader program. It does not claim to settle formal complexity theory in full generality; instead, it internalizes the reduction-correctness lemmas, size-bounded hardness packages, exact finite deciders, counted-search procedures, tractability wrappers, explicit-state upper-bound interfaces, and the scoped witness-over-verifier machinery now used for the stochastic $\textsf{NP}^{\textsf{PP}}$ upper-bound story. In that sense the contribution is closest to a certified reduction-and-decision core for one theorem family rather than to a general-purpose complexity library.

\subsection{Rough Sets, Reducts, and Attribute Reduction}

The closest classical neighbor is rough-set theory and its large literature on reducts and attribute reduction~\cite{pawlak1982rough,jensen2004semantics,jensen2008rough}. That literature asks when a smaller attribute set preserves discernibility or decision-table structure, and it has developed both exact and heuristic methods for reduct computation. There is an obvious family resemblance to the static regime of the present paper: both settings ask which coordinates can be removed without losing decision distinctions.

The difference is not merely terminological. Rough-set reducts are typically formulated relative to indiscernibility relations or decision tables, whereas the object here is the optimizer map of a decision problem and the induced exact preservation of optimal-action sets. That shift matters mathematically as well as computationally. The present paper does not study a single static reduct notion in isolation; it studies one exact certification question and then tracks corresponding predicates across static, stochastic, and sequential regimes. In the present formulation, insufficiency has short counterexample witnesses, sufficiency is universal over such witnesses, and minimum sufficiency collapses to relevance containment in the static regime. The stochastic and sequential layers then introduce conditional-comparison and temporal-structure phenomena that are outside the usual rough-set formulation. The paper's static-collapse theorem, regime hierarchy, and exact-certification trilemma are therefore not direct corollaries of the reduct literature, even though that literature is an important intellectual precursor and deserves explicit credit as a nearby tradition.

The same generosity is due to the broader feature-selection and subset-selection literature in machine learning and combinatorial optimization~\cite{blum1997selection,amaldi1998complexity}. Those works study sparse predictive representations, informative feature subsets, and the computational cost of choosing them. Our setting differs in targeting exact preservation of the optimal-action correspondence rather than prediction loss, classifier complexity, or approximate empirical utility. Still, the shared concern is clear: identifying what information matters and what can be safely omitted.

\subsection{Decision-Theoretic Sufficiency, Informativeness, and Value of Information}

There is also a strong conceptual connection to statistical decision theory and the theory of informativeness. Classical sufficient statistics, value-of-information analysis, and comparison of experiments ask when one information structure is at least as useful for decision making as another~\cite{blackwell1953equivalent,savage1954foundations,raiffa1961applied,howard1966information}. That tradition studies when information can be compressed, compared, or discarded without harming decisions. The optimizer quotient and the exact sufficiency predicate in this paper should be read against that background.

At the same time, the present work studies a different layer of the problem. The focus here is not statistical estimation, asymptotic identifiability, or expected utility under a fixed experiment class. It is the computational certification problem of proving that a coordinate projection preserves the optimal-action correspondence exactly. The connection to classical decision theory is therefore one of ancestry and motivation rather than direct theorem transfer. The paper owes that literature recognition for posing, in older language, the same high-level question: what information is genuinely needed for a decision?

\subsection{Abstraction, Bisimulation, and Homomorphisms in Planning and Reinforcement Learning}

The stochastic and sequential parts of the paper sit near a second major literature: exact abstraction in Markov decision processes, POMDPs, bisimulation-based aggregation, and MDP homomorphisms~\cite{papadimitriou1987mdp,littman1998probplanning,mundhenk2000mdp,givan2003equivalence,ravindran2004algebraic}. This body of work studies when state spaces can be aggregated while preserving value or policy structure, and it has established many of the representation-sensitive complexity phenomena that make richer stochastic and sequential models computationally difficult.

That literature is structurally very close to the present paper and should be credited as such. The main difference is that our predicate is coordinate sufficiency for preserving optimal-action sets under an explicit coordinate-hiding operation. We do not study arbitrary abstract state maps, approximate bisimulation metrics, or value-function approximation guarantees. Instead, we ask for exact certification that a chosen coordinate interface is enough. This exact-certification emphasis is what the complexity classification adds beyond prior abstraction results: the question is not only whether a smaller representation can support good control, but whether one can certify that no omitted coordinate changes the optimal-action correspondence. The contribution here is therefore not abstraction theory in general, but a regime-sensitive complexity theory for one exact decision-relevance predicate tracked coherently across static, stochastic, and sequential settings.

\subsection{Explanations, Anchors, and Minimal Decision Rules}

The paper should also acknowledge adjacent work in explainable AI and local rule-based explanations, especially anchor-style explanations~\cite{ribeiro2018anchors}. Readers will naturally notice the word ``anchor'' and map it to that literature. The resemblance is real: both settings isolate a restricted condition under which a local piece of information suffices to support a decision or prediction. More broadly, explanation methods based on local rules, counterfactuals, or minimal sufficient subsets share the same informal ambition of identifying a small decision-preserving core.

The difference is again exactness and scope. Anchor explanations in XAI are typically local, model-agnostic, approximate, or precision-calibrated. Our anchor-sufficiency problem is an exact certification problem inside a fully specified decision model, and its complexity is governed by an existential choice followed by a universal fiberwise verification condition. So the paper is not an explanation method, but it is in direct conversation with explanation work that asks what small local structure is enough to preserve a decision.

There is an even closer logical neighbor in recent work on reasons behind decisions, prime-implicant style explanations, and exact decision justifications~\cite{darwiche2023reasons}. This is especially worth noting because it shows that exact small decision-preserving structures are a live contemporary topic, not only a classical one. That literature is concerned with exact structural reasons for a decision outcome and therefore lies much nearer to the present paper than generic post-hoc explanation methods do. The overlap is substantial at the conceptual level: both ask for exact small structures that preserve a decision. The difference is that our objects are coordinate sufficiency, minimum sufficiency, and anchor sufficiency inside coordinate-structured decision problems, and the paper's main contribution is a regime-sensitive complexity classification of these predicates rather than a logical semantics of reasons alone.

\subsection{Abduction, Diagnosis, and Exact Explanatory Cores}

The minimum and anchor queries should also be read against the classical literature on abduction and diagnosis~\cite{reiter1987diagnosis,eiter1995abduction}. Reiter-style diagnosis studies which hidden fault sets explain an observation, while logic-based abduction studies the complexity of finding explanatory hypotheses sufficient for a target consequence. Those frameworks share an important structural feature with the present work: a move from verification to explanatory search typically raises complexity because one is no longer merely checking a candidate structure but searching for one.

That kinship is especially relevant for the paper's anchor query. The existential anchor choice is one reason the static anchor problem remains $\Sigma_2^P$-complete even when minimum sufficiency collapses in the static regime. In that sense the anchor problem behaves more like explanatory discovery than like plain verification. The exact predicate is different, because we certify preservation of optimal-action sets under coordinate hiding rather than explanatory entailment of a formula, but the logical-complexity intuition is nearby enough that this literature deserves explicit acknowledgment.

\subsection{Causality, Responsibility, and Structural Explanation}

The paper also belongs in conversation with structural accounts of causality, explanation, and responsibility~\cite{pearl2009causality,spirtes2000causation,halpern2001explanations,chockler2004responsibility}. Those works ask which variables or interventions matter for an outcome, how explanatory structure should be represented, and how responsibility can be assigned to particular factors. This is not the same problem studied here, but it is nearby in spirit: a relevant coordinate is precisely a coordinate whose perturbation can change the decision, and the optimizer quotient isolates exactly the distinctions that matter for optimal-action structure.

The difference, again, is that the present paper focuses on exact decision preservation under coordinate projection and on the computational cost of certifying that preservation. Causality and responsibility frameworks typically emphasize semantic characterization, counterfactual dependence, or intervention-level explanation. The present work instead asks when hidden coordinates can be ignored \emph{without changing the optimizer map}, and how hard it is to certify that claim exactly across static, stochastic, and sequential regimes.

\subsection{Oracle, Query, and Access-Based Lower Bounds}

Query-access lower bounds provide another nearby technique family~\cite{dobzinski2012query}. Our witness-checking duality and access-obstruction results belong to that tradition in spirit. The difference is that the lower bounds are developed around the same fixed sufficiency predicate used throughout the paper, which allows direct comparison between forward evaluation, certification, exact minimization, and restricted-access models without changing the underlying decision relation.

\subsection{Scope and Novelty}

The paper is therefore intentionally synthetic in the good sense. Rough sets and reducts study decision-preserving attribute reduction; feature selection studies informative subsets; statistical decision theory studies informativeness and sufficiency; abstraction and bisimulation literatures study exact aggregation of stochastic and sequential decision processes; explanation methods study local decision-preserving cores; formalized complexity work studies certified reductions and machine-checked decision procedures. The contribution here is to treat exact relevance certification itself as the common object and to develop one coherent complexity-theoretic account around it.

What is new is not merely that these literatures are mentioned together. It is that one predicate, exact preservation of optimal-action distinctions under coordinate hiding, is formalized once and then followed across static, stochastic, and sequential regimes; that the static regime exhibits a non-obvious collapse of minimum sufficiency to relevance containment; that the encoding-sensitive dichotomy isolates a sharp tractable--intractable boundary; that the downstream consequences are pushed through to exact simplification, witness checking, approximation, and the simplicity tax; and that a substantial Lean-backed certification layer is provided for the reduction and finite-decision core. Prior work contains many nearby pieces. The present paper aims to assemble them into a single precise theory of exact relevance certification.
